We build on ReLoRA~\cite{lialin2023stack} and train with successive LoRA cycles.
We freeze pretrained weights $W \in \mathbb{R}^{m \times n}$ and parameterize each cycle update as a low-rank product
\begin{align}
\begin{split}
\delta W &= s\, W_A W_B,\\
W_A &\in \mathbb{R}^{in \times r},\qquad
W_B \in \mathbb{R}^{r \times out},
\end{split}
\label{eq:1}
\end{align}
where $s \in \mathbb{R}$ is a fixed scaling factor (typically $\alpha_{\mathrm{LoRA}}/r$).

Every $q$ steps we merge the current adapter into the base weights and start a new LoRA cycle.
After $N$ cycles the total update is
\begin{align}
\begin{split}
\Delta W
&= \sum_{i=1}^{N} \delta W_i
= \sum_{i=1}^{N} s\, W_A^{(i)} W_B^{(i)}.
\end{split}
\label{eq:2}
\end{align}
If successive factors reuse the same input subspace, the sum in~\eqref{eq:2} need not increase rank.
ReLoRA relies on random reinitialization and optimizer pruning to encourage new directions, but does not enforce independence of subspaces.

\paragraph{Column-space orthogonality at reinit.}
Our modification is to force each new cycle to start in an input subspace orthogonal to previous ones.
Let $U$ be an orthonormal basis of the column space spanned by all previously merged adapters $\{W_A^{(i)}\}_{i<k}$.
At the beginning of cycle $k$ we reinitialize $W_A^{(k)}, W_B^{(k)}$ as in ReLoRA and then project
\begin{equation}
W_A^{(k)} \leftarrow \bigl(I - U U^\top\bigr) W_A^{(k)}.
\label{eq:orth}
\end{equation}
Thus $W_A^{(k)\top} U = 0$ at cycle start: the new adapter uses input directions disjoint from past cycles.
We update $U$ after each merge by appending an orthonormal basis of the just-merged $W_A$ (e.g., via QR).
Orthogonality is enforced \emph{only at reinit}; during the cycle the factors may drift under gradient updates.
This yields a multi-cycle LoRA scheme whose successive updates are initialized to be column-space orthogonal, aiming for
$\mathrm{rank}(\Delta W)$ to grow with the number of cycles.

\paragraph{Transition momentum.}
A hard full merge assumes that the current cycle has already absorbed a useful rank-$r$ update into $W$.
In practice merges are frequent, so $W_A,W_B$ may be far from converged: discarding them and forcing a purely orthogonal $W_A$ can throw away unfinished work and induce a suboptimal transition into the next cycle.
We therefore introduce a momentum coefficient $\mu \in [0,1]$ that softens the merge-and-reinit step.
Writing $W_A^{\mathrm{old}},W_B^{\mathrm{old}}$ for the adapters just before the merge and $\delta W = s\, W_A^{\mathrm{old}} W_B^{\mathrm{old}}$, we perform a \emph{partial} merge
\begin{equation}
W \leftarrow W + (1-\mu)\, \delta W,
\label{eq:partial_merge}
\end{equation}
reinitialize fresh factors $W_A^{\mathrm{new}},W_B^{\mathrm{new}}$ (with~\eqref{eq:orth} when using SuperReLoRa), and blend
\begin{align}
\begin{split}
W_A &\leftarrow (1-\mu)\, W_A^{\mathrm{new}} + \mu\, W_A^{\mathrm{old}},\\
W_B &\leftarrow (1-\mu)\, W_B^{\mathrm{new}} + \mu\, W_B^{\mathrm{old}}.
\end{split}
\label{eq:reinit_momentum}
\end{align}
Thus $\mu{=}0$ recovers the classical hard reset, while $\mu{>}0$ keeps a fraction of the previous cycle in both the base weights (via the residual adapter) and the new factors.
Adam moments are pruned as in ReLoRA and then scaled by $\mu$ so that optimizer momentum is carried across the boundary rather than fully discarded.
As in ReLoRA, we also restart the learning rate with a jagged warmup schedule after each merge.
The full procedure is given in~\Algref{alg:superrelora}.

\begin{algorithm}[H]
\caption{SuperReLoRa with transition momentum.
\(\theta\) denotes model parameters; \(\hat{\theta}\) the model after replacing linear layers by LoRA adapters; \(M,V\) are Adam states; \(\eta\) follows a jagged schedule; \(q\) is the merge/reinit period; \(\mu\) is the transition momentum; \(U\) stores an orthonormal basis of past adapter column spaces.}
\label{alg:superrelora}
\begin{algorithmic}[1]
\Require \(\theta\), \(M\), \(V\), \(\eta\), \(q\), \(\mu\)
\State \(U \leftarrow [\,]\) \{empty basis\}
\For{\(t\) in warm-start steps}
    \State Update \(\theta\), \(M\), \(V\), \(\eta\) \{standard training\}
\EndFor
\For{layer in model layers}
    \If{layer is linear}
        \State Replace layer by LoRA\((W, W_A, W_B)\); freeze \(W\)
    \EndIf
\EndFor
\For{\(t\) in training steps}
    \State Update \(\hat{\theta}\), \(M\), \(V\) \{LoRA training step\}
    \If{\(\mathrm{MOD}(t,q)=0\)}
        \For{\(l\) in linear layers}
            \State \(\delta W \leftarrow s\, W_A W_B\)
            \State \(W \leftarrow W + (1-\mu)\, \delta W\) \{partial merge\}
            \State Append orthonormal basis of \(\mathrm{col}(W_A)\) to \(U\); re-orthonormalize \(U\)
            \State \(W_A^{\mathrm{old}} \leftarrow W_A\); \(W_B^{\mathrm{old}} \leftarrow W_B\)
            \State \(W_A \leftarrow \mathrm{kaiming\_init}(W_A)\); \(W_B \leftarrow 0\) \{fresh reinit\}
            \State \(W_A \leftarrow (I - U U^\top) W_A\) \{column-space orthogonalization\}
            \State \(W_A \leftarrow (1-\mu)\, W_A + \mu\, W_A^{\mathrm{old}}\) \{momentum blend\}
            \State \(W_B \leftarrow (1-\mu)\, W_B + \mu\, W_B^{\mathrm{old}}\)
            \State \(M_{W_A},V_{W_A},M_{W_B},V_{W_B} \leftarrow \mathrm{prune}(\cdot)\); then scale by \(\mu\)
        \EndFor
        \State Restart \(\eta\) warmup
    \EndIf
\EndFor
\State \textbf{return} \(\theta\)
\end{algorithmic}
\end{algorithm}
